\subsection{Simulación XOR con Detección de Ciclos}

\graybold{Preguntas guía} - ¿cuándo usarlo?

\begin{itemize}
\item ¿El estado del sistema se puede representar como un vector de bits, y cada operación disponible es un XOR con una máscara fija?
\item ¿Las operaciones se aplican en un orden CÍCLICO fijo y predecible (no libre elección)?
\item ¿Necesito detectar si el proceso entra en un ciclo infinito sin nunca alcanzar el estado objetivo?
\end{itemize}

\graybold{¿Para qué sirve?} Simular un proceso determinista donde el estado (aquí, qué bombillas están encendidas) cambia por XOR con máscaras fijas aplicadas en orden cíclico, deteniéndose al llegar al estado objetivo o detectando que nunca se alcanzará. Es aplicable a cualquier proceso de estado finito con transiciones deterministas.

\graybold{Complejidad:} \bigO{N} en la práctica (ver fundamento teórico), a pesar de que el espacio de estados aparenta ser exponencial en M.

\graybold{Fundamento teórico:} Presionar los N interruptores en orden fijo $1..N$ corresponde a aplicar sucesivamente XOR con N máscaras fijas. Después de un CICLO COMPLETO (los N interruptores), el efecto neto sobre el estado es XOR con una constante fija $C$ = XOR de todas las máscaras de los interruptores, sin importar el estado en que se esté. Como $v \oplus C \oplus C = v$, el estado al INICIO de cada ciclo solo puede tomar 2 valores posibles (el original, o el original $\oplus\, C$), alternando entre ambos. Esto acota el número de estados distintos alcanzables — y por lo tanto el número de pasos antes de repetir un estado o llegar a 0 — a, como mucho, $2N$: por eso la detección de ciclos con un conjunto de estados visitados nunca necesita más de un puñado de iteraciones, sin importar cuán grande sea M.

\graybold{Ejercicio aplicado}

Un conserje presiona N interruptores en orden cíclico $1,2,\ldots,N,1,2,\ldots$ Cada interruptor alterna un subconjunto fijo de M bombillas. Dado el estado inicial de las bombillas encendidas, determinar cuántas veces hay que presionar interruptores hasta que TODAS las bombillas queden apagadas simultáneamente, o -1 si nunca ocurre.

\graybold{I/O:} N,M ≤ 1000 → readline por línea (patrón 1) es más que suficiente, dado que el propio algoritmo termina en a lo más \textasciitilde{}2N iteraciones (verificado: <0.03s incluso en el peor caso N=M=1000).

\graybold{Código solución}

\begin{lstlisting}[language=Python]
import sys
input = sys.stdin.buffer.readline

n, m = map(int, input().split())
ldata = list(map(int, input().split()))   # ldata[0]=cantidad, resto=bombillas encendidas

switch = []
for i in range(n):
    sdata = list(map(int, input().split()))
    tmp = 0
    for j in range(1, sdata[0] + 1):
        aux = 1 << (m - sdata[j])          # bit que representa esa bombilla
        tmp += aux
    switch.append(tmp)                     # mascara XOR de este interruptor

bulbs = 0
for j in range(1, ldata[0] + 1):
    aux = 1 << (m - ldata[j])
    bulbs += aux                           # estado inicial como bitmask

idx = 0
res = 0
seen = set()                               # detecta si (estado, interruptor actual) se repite
while True:
    if bulbs == 0:
        break                              # todas las bombillas apagadas: listo
    if (bulbs, idx) in seen:
        res = -1                           # ya visitamos este estado exacto: nunca terminara
        break
    seen.add((bulbs, idx))
    info = switch[idx]

    bulbs ^= info                          # aplica el interruptor actual

    if idx + 1 == n:
        idx = 0                            # vuelve al primer interruptor (orden ciclico)
    else:
        idx += 1
    res += 1

sys.stdout.write(str(res))
\end{lstlisting}